% =====================================================================
% The Box Universe -- draft v0.1 (2026-07-29, dispatch 43)
% Written per docs/dev/HANDOFF_PAPER2_WRITE.md v1.1 (sole authoritative
% writing brief; approved skeleton HANDOFF_PAPER_BOX v0.3).
% Status: full prose draft; figures 1-8 are planned placeholders
%         (machine generation wave pending -- Sec. VIII table maps every
%         numeric claim to its verification gate).
% Input canon: docs/BOX_UNIVERSE.md v1.6, docs/PHYSICS.md v1.14,
%              docs/DERIVATIONS.md Sec. 17, beta v1.34 (V29 PASS).
% Title: option (c), adopted by author decision (dispatch 44,
%        2026-07-29).
% v0.3: the 0.92-vs-a relation stated exactly (probe-agreement epoch
%        a=2D/Kt~1.85; a=2 point value = anchor value identically;
%        numeric coincidence 2^{4/3}~e^{0.9233} to 0.10%).
% v0.2: title (c) adopted; D/K_t calibration given its exact closed
%        form (2/3)ln4 with the strong-field caveat, incorporating the
%        external reviewer's independent computation (2026-07-26).
% Build:  pdflatex dfm-paper2.tex  (RevTeX 4.2, as paper 1)
% =====================================================================
\documentclass[aps,reprint,amsmath,amssymb,nofootinbib,superscriptaddress,floatfix]{revtex4-2}

\usepackage{graphicx}
\graphicspath{{figures/}}
\usepackage{bm}
\usepackage{hyperref}
\hypersetup{
  pdftitle={The Box Universe: Expansion, Time, and Light in a Relational
    Toy Model},
  pdfsubject={Educational toy model; expansion; Hubble tension analogy}}

\begin{document}

\title{The Box Universe: Expansion, Time, and Light in a Relational Toy Model}

\author{Tetsuya Imamura}
\affiliation{Independent researcher, Japan}

\date{2026-07-29 (draft v0.3)}

\begin{abstract}
We extend Determinacy-Field Mechanics (DFM), a two-dimensional
relational toy model introduced in a companion paper, with a
\emph{box universe}: the approximation that everything outside a
finite region is replaced by a single simplified background.
Within this one construction we demonstrate three things.
(i)~The box makes the classic relational questions computable:
uniform translation of the box is undetectable from inside
(exactly), while rotation and expansion couple to the interior
through a measured gain factor $g(r)$ with $g(0)=1/2$ for a uniform
ring, matching the analytic average $\langle\hat n(\hat n\cdot
\bm r)\rangle=\bm r/2$.
(ii)~Expansion appears in three separable layers---prescribed
kinematics, wavelength semantics, and solved dynamics---each
machine-verified; a pressure-driven free box accelerates
($a_{\mathrm{eff}}=2.18$) while its gravity-only control
recollapses.
(iii)~Most importantly, the model exhibits \emph{probe-dependent
expansion measurement}: a geometric probe and a velocity-decay probe
applied to the same box disagree by the analytic factor
$H_w/H_{\mathrm{geo}}=2(D/K_t)\,a^{-d_P}$, verified numerically to
$5.3\times10^{-4}$.
With the single calibration $D/K_t=0.92$ the disagreement at $a=2$
equals $0.92$, the same magnitude as the ratio of the two measured
values of the Hubble constant ($67.4/73.0\approx0.923$).
We present this as a \emph{structural analogy}---a minimal mechanism
by which the measured expansion rate depends on the probe---not as an
explanation of the Hubble tension.
A final section assembles the model's optically dark gravitating
bodies and lists which phenomenological requirements of dark matter
they meet inside the model, together with the measured limits of that
construction.
All numeric claims are tied to automated verification gates that run
in continuous integration.
\end{abstract}

\maketitle

% ---------------------------------------------------------------------
\section{Introduction}
\label{sec:intro}

The companion paper introduced Determinacy-Field Mechanics (DFM), an
educational toy model in which space is purely relational: every
mechanical property of a point is set by the surrounding masses
through a scalar \emph{determinacy field} $W$, and inertia, gravity,
clock rate, and light propagation are all read from it.
That paper studied local consequences---frame dragging, rotation
curves, clocks, and lensing---on an effectively infinite background.

This paper studies the opposite end of the approximation: what happens
when ``the rest of the universe'' is deliberately replaced by a single
simplified object, which we call the \emph{box}.
The box universe is not a new axiom; it is the statement that beyond a
chosen radius the determinacy sources are summarized by a uniform
background term $D_0$, a box determinacy $D$ with scale factor $a(t)$,
and optionally a uniform gravity vector.
Every cosmological question the toy model can ask---What expands?
What is comoving? What does a redshift mean? Can expansion be
\emph{solved} rather than prescribed?---becomes a question about the
box, and, because the model is small and deterministic, each answer
can be machine-verified.

Our position is unchanged from the companion paper: DFM is an
educational \emph{computational thought experiment}, not a proposal
about nature (Sec.~\ref{sec:negative}).
The interest of the box universe is that it makes textbook
distinctions---kinematic expansion versus dynamic expansion,
coordinate recession versus measured redshift, probe versus
quantity---concrete enough for a student to toggle in a browser, with
a continuous-integration suite asserting every number in this paper.

The paper is organized as follows.
Section~\ref{sec:box} defines the box and establishes what the inside
can and cannot detect.
Section~\ref{sec:layers} separates expansion into three verified
layers.
Section~\ref{sec:probe} contains the main result: probe-dependent
expansion measurement and its structural analogy with the Hubble
tension.
Section~\ref{sec:dark} assembles the model's optically dark bodies.
Section~\ref{sec:negative} states six negative claims that are part of
the thesis.
Section~\ref{sec:repro} maps every claim to its verification gate.

% ---------------------------------------------------------------------
\section{The box as an approximation}
\label{sec:box}

\subsection{Semantics}
\label{sec:semantics}

The box universe is defined by six semantic commitments (canonical
statement in the project's box-universe document, Sec.~12):
(1)~outside a finite region, determinacy sources are summarized, not
simulated;
(2)~the summary consists of a uniform background $D_0$, a box term
$W_B=D\,a^{-d_P}$ carried by a scale factor $a(t)$, and optionally a
uniform gravity vector---conceptually these are all \emph{box
parameters};
(3)~interior physics is unchanged: the same axioms read the total
field $W$ wherever it comes from;
(4)~the box may be \emph{prescribed} (a regulated $a(t)$) or
\emph{free} (walls made of ordinary masses whose motion solves
$a(t)$);
(5)~comoving coordinates and wavelength semantics are definitions
layered on top, each with its own verification; and
(6)~no statement about the real universe is implied by any of it.

\subsection{What the inside cannot detect}
\label{sec:t10}

Uniform translation of the entire box is exactly undetectable from
inside: the relational axioms depend only on relative configuration,
and the verification suite asserts bit-identical interior evolution
under a boosted box.
Rotation and expansion are different: they enter through the
determinacy frame with a measured \emph{gain} $g(r)$.
For a uniform ring of wall masses the analytic average of the
direction-projected position over sources gives
\begin{equation}
  \bigl\langle \hat n\,(\hat n\cdot\bm r)\bigr\rangle=\tfrac{1}{2}\bm r ,
  \label{eq:gain}
\end{equation}
so the interior center couples to wall rotation and expansion with
gain $g(0)=1/2$; the measured value agrees (gates V23a/V24a).
This is the toy-model version of the classic relational statement:
absolute uniform motion is meaningless, but absolute rotation and
expansion of ``everything else'' are interior-observable through the
frame they drag.
[Fig.~1: box hierarchy schematic; Fig.~2: $g(r)$ measured vs.\
analytic --- planned.]

% ---------------------------------------------------------------------
\section{Three layers of expansion}
\label{sec:layers}

\subsection{Layer 1: prescribed kinematics}

With a regulated $a(t)$ the box provides FLRW-style bookkeeping
without dynamics.
The comoving coordinate $\chi=r/a$ of a fully comoving particle is
conserved to $1.5\times10^{-5}$ over the verification run (V24b).
Bound systems follow the box only fractionally: the measured orbit
scaling $d\propto a^{2\varphi_B}$ interpolates with the determinacy
share $\varphi_B$ of the box at the orbit (V25: net exponent $1.61$
at $\varphi_B{\to}1$ and $0.00$ at $\varphi_B{\to}0$), giving the
textbook ``bound systems do not expand'' as the $\varphi_B\to 0$ end
of a measured dial.

\subsection{Layer 2: wavelength semantics}

Redshift in the box is a \emph{semantic} rule: emitted light retains
the wavelength ratio $1+z=a_{\mathrm{obs}}/a_{\mathrm{emit}}$.
The suite emits three photons (A/B/C) at different epochs of the same
run and checks all three simultaneously (gate \texttt{box.photon-abc},
bookkeeping closed below $10^{-9}$).
This is a definition layered on the kinematics, not a derivation from
photon dynamics (negative claim 2 of this paper); its value is that
it reproduces, at the level of the \emph{equations}, the observed
consistency relations---supernova time dilation $\propto(1+z)$ and
$T_{\mathrm{CMB}}(z)=T_0(1+z)$---so a student can see exactly which
part of those relations is semantics and which is physics.

\subsection{Layer 3: solved dynamics --- the free box}

Replacing the regulated box by 64 free wall masses turns $a(t)$ into
an output.
With spin pressure on, the measured effective scale factor
accelerates to $a_{\mathrm{eff}}=2.18$ (the pair repulsion decays as
$1/d$, slower than gravity, so it wins at large separation); with
pressure removed ($k_{\mathrm{Rep}}=0$) the same box coasts to
$a_{\mathrm{eff}}\approx1.005$ and recollapses.
The run is a fully closed system: reservoir bookkeeping is
identically zero and momentum/angular-momentum drift is below
$4\times10^{-8}$ (gates \texttt{freebox.*}).
This is a mechanism demonstration in the spirit of Friedmann
cosmology---attraction decelerates, pressure-like repulsion can
accelerate---with no claim of reproducing the Friedmann equations
(negative claim 5).
[Fig.~5: $a_{\mathrm{eff}}(t)$ with and without pressure --- planned.]

\subsection{Aside: an arrow of time in a reversible box}

With the leapfrog integrator the model is exactly reversible: flipping
all velocities and spins returns the initial state to RMS
$1.3\times10^{-3}$.
Turning on one dissipative channel ($\mu_F=0.3$) breaks the return
(RMS $10.8$); the default semi-implicit integrator, though
deterministic, fails to rewind by itself (RMS $0.92$, a $709\times$
degradation).
The arrow of time in the box is thus exhibited as a property of
dissipation, not of the dynamics' determinism (gates
\texttt{echo.*}).
This section is deliberately brief; a fuller treatment is future
work.

% ---------------------------------------------------------------------
\section{Probe-dependent expansion measurement}
\label{sec:probe}

\subsection{Hubble anti-friction and the ratio standard}

In DFM the peculiar velocity $w$ of a free particle is conserved in
the absolute bookkeeping of the model even while the box expands
(measured drift $<10^{-3}$; gates V27/V28).
What decays is the \emph{ratio} $R=w/c_{\mathrm{loc}}$: the local
light speed rises as the box determinacy dilutes,
$c_{\mathrm{loc}}=c_0\,e^{-2\psi}$ with $\psi=W_B/K_t$, giving
\begin{equation}
  R(a)=\exp\!\Bigl[-\tfrac{2D}{K_t}\bigl(1-a^{-d_P}\bigr)\Bigr],
  \label{eq:R}
\end{equation}
verified to $2.5\times10^{-4}$ at $a=2$ (V28).
Measured against light---the only standard an interior observer
has---free motion \emph{does} decay, in the same direction as the
FLRW $w\propto1/a$ friction, although with a different functional
form.

\subsection{Two probes, one box}
\label{sec:v29}

Now let two observers measure the expansion rate of the same box.
The \emph{geometric} observer tracks comoving structure and reports
$H_{\mathrm{geo}}=\Delta\ln a/\Delta t$.
The \emph{velocity} observer assumes FLRW kinematics
($w/c\propto1/a$) and infers an expansion rate from the measured
decay of $R$:
$H_w=-\Delta\ln R/\Delta t$.
From Eq.~(\ref{eq:R}),
\begin{equation}
  \frac{H_w}{H_{\mathrm{geo}}}=\frac{2D}{K_t}\,d_P\,a^{-d_P},
  \label{eq:ratio}
\end{equation}
which depends on the epoch $a$ and on the box calibration $D/K_t$.
The two probes disagree, not because either is wrong, but because the
model's decay law is not the FLRW one the second observer assumed.

Verification gate V29 measures both probes in one run over three
windows ($a\colon1\to1.5\to2\to3$, $d_P=1$) and compares the measured
ratio with the log-window average of Eq.~(\ref{eq:ratio}); the
largest relative error is $5.3\times10^{-4}$ against a tolerance of
$10^{-2}$, with $w$ conserved to $4.8\times10^{-4}$.
[Fig.~6: $H_w/H_{\mathrm{geo}}$ vs.\ $a$ for several $D/K_t$ ---
planned; the underlying scan is specified in the project's experiment
design 4-66.]

\subsection{A structural analogy with the Hubble tension}
\label{sec:tension}

The two accepted measurements of the real Hubble constant disagree:
$H_0=67.4$ (CMB fit) versus $73.0$ (local distance ladder),
a ratio of $0.923$.
In the box universe the calibration has an exact closed form: the
choice that makes the ratio standard $R$ of Eq.~(\ref{eq:R}) fall to
exactly $1/4$ at $a=4$ is
\begin{equation}
  \frac{D}{K_t}=\frac{\ln 4}{2\,(1-1/4)}=\tfrac{2}{3}\ln 4
  \approx 0.9242 ,
\end{equation}
first computed independently by an external review of the project
(2026-07-26), and with it the probe ratio at $a=2$ is
\begin{equation}
  \frac{H_w}{H_{\mathrm{geo}}}\Big|_{a=2}=\frac{2D}{K_t}\cdot\frac12
  =\frac{D}{K_t}\approx 0.92 ,
\end{equation}
the same magnitude as the observed tension ratio.
The relation between the two appearances of $0.92$ is worth stating
exactly (dispatch-44/45 audit): the probe ratio of
Eq.~(\ref{eq:ratio}) equals $D/K_t$ at $a=2$ \emph{identically}, for
any calibration, so ``$0.92$ at $a=2$'' is simply the $a{=}4$ anchor
value read at $a=2$; the probes agree ($H_w=H_{\mathrm{geo}}$) at
$a=2D/K_t\approx1.85$; and $R$ crosses the FLRW curve $1/a$ exactly
twice, at $a=1$ and (by construction) $a=4$, decaying faster in
between.
The numeric agreement with the tension ratio reduces to the identity
$2^{4/3}\approx e^{0.9233}$, true to $0.10\%$---a coincidence of
constants, which is why negative claim~6 treats it as calibration,
not evidence.
One honesty note accompanies the calibration: $\psi=D/K_t\approx0.92$
is a \emph{strong-field} setting of the toy
($c_{\mathrm{loc}}/c_0=e^{-2\psi}\approx0.16$), far outside the
weak-field regime $|\psi|\ll1$ in which the companion paper's GR
calibrations live; nothing in this section inherits those
calibrations.
We state precisely what this is and is not.
It \emph{is} a minimal, fully inspectable mechanism in which the
measured expansion rate depends on the probe---early-universe-style
global fits and late-universe-style local kinematics need not agree
if the underlying decay law differs from the assumed one.
It is \emph{not} an explanation of the Hubble tension: the model is
two-dimensional, non-Lorentzian, and the coincidence of $0.92$ with
$0.923$ is a calibration, not a prediction (negative claims 3 and 6).

The functional form is the entire content of the effect.
If the model's light-speed law were replaced by the exact inverse
proportionality $c_{\mathrm{loc}}\propto1/W$ (ledger item 4-66), then
$R\propto a^{-d_P}$ exactly, both probes would coincide identically
($H_w/H_{\mathrm{geo}}\equiv d_P$), and the probe dependence would
vanish.
The comparison experiment that separates the two forms---an
$a\times D/K_t$ scan of Eq.~(\ref{eq:ratio}) against the constant
line---is specified as a design document in the repository and doubles
as the data source for Fig.~6.
The choice between the two laws is thus \emph{observable from inside
the model}, which is the pedagogical point: what looks like a
measurement discrepancy can be a statement about which decay law
nature (here: the axiom set) actually uses.

Because $\rho(a)\equiv2(D/K_t)a^{-d_P}$ falls with $a$, the
disagreement is largest early and vanishes late, giving the toy model
an early/late structure qualitatively reminiscent of the real
tension's framing---again as structure, not as fit.

% ---------------------------------------------------------------------
\section{Optically dark gravitating bodies}
\label{sec:dark}

The companion paper showed that DFM's rotating cores enhance outer
rotation curves through frame dragging (measured outer boost
$1.332\pm0.009$ over eight seeds).
Here we assemble the model's \emph{darkness} mechanisms and ask,
inside the model, which phenomenological requirements of dark matter
a single object can meet.

\subsection{Two independent darkness mechanisms}

\emph{Self-light dimming} (parameter \texttt{lightSweep}) attenuates
the light a body emits; since v1.33 it can be computed from the spin
state as $l_{S,\mathrm{eff}}=\min(1,|s|R/c_{\mathrm{surf}})$ rather
than set by hand, and for the model's black-hole-like preset the
mechanistic value is $0.56$ against the toy value $1$---a gap the
interface exposes rather than hides.
\emph{External-ray sweep-out} is a strong-field effect of a rapidly
spinning core on \emph{incoming} rays: non-escape rises over
$s=0.5$--$1.0$, and the effect vanishes entirely under the physical
lock $K_t=c^2/G$ (measured; the mechanism lives in the toy's strong
regime only).
The two mechanisms are machine-verified to be orthogonal: ray
endpoints agree to all digits as \texttt{lightSweep} is scanned from
0 to 1 (project derivations, Sec.~17).
[Fig.~7: ray sweep at zero and high spin --- planned.]

\subsection{Requirements met, and limits}

Inside the model, a heavy, fast-spinning, optically dark core
(i)~sources $W$ and hence gravity; (ii)~is dark by two independent
mechanisms; and (iii)~contributes to outer rotation-curve enhancement
through dragging, with the radius-resolved profile showing the boost
concentrated inward ($v/v_{\mathrm{bar}}$ enhancement from $1.58$ at
inner radii to $1.15$ outside, eight-seed statistics).
It therefore meets, \emph{inside the model}, the phenomenological
requirements ``dark but gravitating, and relevant to rotation
curves''.
The limits are measured with the same care: the inverse-design
problem (choose parameters to hit a prescribed outer velocity) is
accurate as a forward predictor (median residual ${\approx}2\%$) but
saturates in azimuthal velocity before reaching the outermost bands,
and earlier claims that dragging alone sustains a specific galaxy
preset's arms were withdrawn when knockout tests showed a composite
pathway.
None of this licenses any claim about real dark matter (negative
claim 4); the value is that requirement-level talk about darkness
becomes an explicit, falsifiable checklist in a system small enough
to audit.
[Fig.~8: $v/v_{\mathrm{bar}}(r)$ --- planned.]

% ---------------------------------------------------------------------
\section{Limitations and negative claims}
\label{sec:negative}

As in the companion paper, these are part of the thesis.

\subsection{Negative claim 1}
The box universe is not a description of the real universe: the model
is two-dimensional, has no Lorentz invariance, and no causal
propagation of the fields.

\subsection{Negative claim 2}
The wavelength-retention rule is semantics, not a derivation from
photon dynamics.

\subsection{Negative claim 3}
Probe-dependent $H$ is presented as a structural analogy; no claim is
made that it explains the Hubble tension.

\subsection{Negative claim 4}
The optically dark bodies of Sec.~\ref{sec:dark} are not proposed as
candidates for real dark matter; they meet requirements \emph{inside
the model} only.

\subsection{Negative claim 5}
The free box demonstrates a mechanism (attraction decelerates,
pressure accelerates); it does not reproduce the Friedmann equations.

\subsection{Negative claim 6}
$D/K_t=0.92$ is a calibration chosen to match the tension ratio's
magnitude; the numeric agreement is not evidence for anything.

% ---------------------------------------------------------------------
\section{Conclusion}
\label{sec:conclusion}

A single deliberately crude approximation---replace everything far
away by a box---is enough to make the standard cosmological
distinctions computable and auditable: what is undetectable
(translation) versus gain-coupled (rotation, expansion); what is
prescribed versus semantic versus solved about expansion; and, at the
center of this paper, the fact that \emph{the measured expansion rate
is a property of the probe as much as of the box}, with the entire
effect carried by the functional form of one decay law.
Every number in this paper is asserted by an automated gate
(Sec.~\ref{sec:repro}), and the model runs in a browser, so each
claim can be re-derived, re-measured, or broken by a student in an
afternoon.
Future work: machine-generated figures for this draft; the
$a\times D/K_t$ comparison scan (design complete); an optional light
probe $H_{\mathrm{light}}$ alongside the two probes of
Sec.~\ref{sec:v29}; and a fuller treatment of the reversibility
demonstration.

% ---------------------------------------------------------------------
\section{Verification map}
\label{sec:repro}

Every numeric claim above maps to an automated verification gate in
the public repository (gates run on every commit for both the release
and beta targets):

\begin{center}
\begin{tabular}{ll}
Claim & Gate \\
\hline
$g(0)=1/2$ (Eq.~\ref{eq:gain}) & V23a / V24a \\
$\chi=r/a$ conserved ($1.5\times10^{-5}$) & V24b \\
$d\propto a^{2\varphi_B}$ ($1.61/0.00$) & V25 \\
prescribed power expansion & V26 \\
FLRW friction contrast ($0.03\%$) & V27 \\
$R(a)$ decay (Eq.~\ref{eq:R}, $2.5\times10^{-4}$) & V28 \\
probe ratio (Eq.~\ref{eq:ratio}, $5.3\times10^{-4}$) & V29 \\
A/B/C wavelengths ($<10^{-9}$) & \texttt{box.photon-abc} \\
free box ($a_{\mathrm{eff}}=2.18$; closed books) & \texttt{freebox.*} \\
arrow of time ($1.3\times10^{-3}$ vs.\ $10.8$) & \texttt{echo.*} \\
darkness orthogonality; sweep-out & derivations Sec.~17 \\
outer boost $1.332\pm0.009$; profile & seeds statistics \\
claim--text consistency (33 claims) & \texttt{claims.sync} \\
\end{tabular}
\end{center}

\end{document}
